\subsection{The distance operator and its spectral subspaces}\label{subsec:spectral-distance-operator}
\paraid{p-834e7a7b7a65}
Fix a nonempty compact metric space
$\MetricSpace{\BaseCarrier}{\MetricSymbol}$.
Equip the space
\[
 \HilbertDomain _\BaseCarrier =\LebesgueSpace^2(\BaseCarrier ,\SelectedMeasure{\BaseCarrier}{\MetricSymbol})
\]
with the inner product
\[
 \langle \HilbertFunction ,\HilbertTestFunction \rangle_{\HilbertDomain _\BaseCarrier }=\int_\BaseCarrier  \HilbertFunction  \HilbertTestFunction \,\IntegrationDifferential \SelectedMeasure{\BaseCarrier}{\MetricSymbol}.
\]
This is a separable real Hilbert space.
Here separability follows from density of continuous functions in
$\LebesgueSpace^2$
and the separability of
$\ContinuousFunctions(\BaseCarrier )$
for compact metric spaces
\cite[Section~6.4]{Muscat2024}.
We define the \emph{distance operator}
$\DistanceOperator{\BaseCarrier}{\MetricSymbol}\colon\HilbertDomain_\BaseCarrier\to\HilbertDomain_\BaseCarrier$
as follows.
For every
$\HilbertFunction\in\HilbertDomain_\BaseCarrier$
and
$\BasePoint\in\BaseCarrier$,
set
\[
 \DistanceOperator{\BaseCarrier}{\MetricSymbol}\HilbertFunction (\BasePoint )
 =\int_\BaseCarrier  \MetricSymbol (\BasePoint ,\IntegrationPoint )\HilbertFunction (\IntegrationPoint )\,\IntegrationDifferential \SelectedMeasure{\BaseCarrier}{\MetricSymbol}(\IntegrationPoint ).
\]
The integral is independent of the representative of
$\HilbertFunction $.
This is the distance kernel operator of
\cite[Definition~11]{MariaOudotSolomon2020}.

\begin{lemma}\label{lem:distance-operator}
\paraid{p-fd7cf639024a}
The operator
$\DistanceOperator{\BaseCarrier}{\MetricSymbol}\colon \HilbertDomain _\BaseCarrier \to \HilbertDomain _\BaseCarrier $
is compact and self-adjoint.
It maps
$\HilbertDomain _\BaseCarrier $
into
$\ContinuousFunctions(\BaseCarrier )$,
and for every
$\HilbertFunction \in \HilbertDomain _\BaseCarrier $
and
$\BasePoint ,\ComparisonPoint \in \BaseCarrier $,
\begin{equation}\label{eq:operator-uniform-bound-1}
 \norm{\DistanceOperator{\BaseCarrier}{\MetricSymbol}\HilbertFunction }_\infty
 \leq\diam\MetricSpace{\BaseCarrier }{\MetricSymbol }\norm{\HilbertFunction }_2,
\end{equation}
and
\begin{equation}\label{eq:operator-uniform-bound-2}
 |\DistanceOperator{\BaseCarrier}{\MetricSymbol}\HilbertFunction (\BasePoint )-\DistanceOperator{\BaseCarrier}{\MetricSymbol}\HilbertFunction (\ComparisonPoint )|
 \leq \MetricSymbol (\BasePoint ,\ComparisonPoint )\norm{\HilbertFunction }_2.
\end{equation}
Every eigenfunction with nonzero eigenvalue has a unique continuous
representative.
\end{lemma}

\paraid{p-e83149336582}
Compactness and self-adjointness of the distance operator are also proved in
\cite[Propositions~12 and~13]{MariaOudotSolomon2020}.

\begin{proof}
\paraid{p-3eab3569fc53}
Let
$\HilbertFunction\in\HilbertDomain_\BaseCarrier$
and
$\BasePoint,\ComparisonPoint\in\BaseCarrier$.
By the Cauchy--Schwarz inequality
\cite[Lemma~3.2.27]{BreitGmeineder2026}
and
$\SelectedMeasure{\BaseCarrier}{\MetricSymbol}(\BaseCarrier)=1$,
we obtain
\[
 \int_\BaseCarrier |\HilbertFunction |\,\IntegrationDifferential \SelectedMeasure{\BaseCarrier}{\MetricSymbol}\leq\norm{\HilbertFunction }_2.
\]
Consequently,
\begin{equation}\label{eq:distance-operator-pointwise-bounds-1}
|\DistanceOperator{\BaseCarrier}{\MetricSymbol}\HilbertFunction (\BasePoint )|
 \leq\int_\BaseCarrier  \MetricSymbol (\BasePoint ,\IntegrationPoint )|\HilbertFunction (\IntegrationPoint )|\,\IntegrationDifferential \SelectedMeasure{\BaseCarrier}{\MetricSymbol}(\IntegrationPoint )
 \leq\diam\MetricSpace{\BaseCarrier }{\MetricSymbol }\norm{\HilbertFunction }_2,
\end{equation}
and
\begin{equation}\label{eq:distance-operator-pointwise-bounds-2}
\begin{aligned}
|\DistanceOperator{\BaseCarrier}{\MetricSymbol}\HilbertFunction (\BasePoint )-\DistanceOperator{\BaseCarrier}{\MetricSymbol}\HilbertFunction (\ComparisonPoint )|
 &\leq\int_\BaseCarrier |\MetricSymbol (\BasePoint ,\IntegrationPoint )-\MetricSymbol (\ComparisonPoint ,\IntegrationPoint )|\,|\HilbertFunction (\IntegrationPoint )|\,\IntegrationDifferential \SelectedMeasure{\BaseCarrier}{\MetricSymbol}(\IntegrationPoint )\\
 &\leq \MetricSymbol (\BasePoint ,\ComparisonPoint )\norm{\HilbertFunction }_2.
\end{aligned}
\end{equation}
The estimates \eqref{eq:distance-operator-pointwise-bounds-1} and
\eqref{eq:distance-operator-pointwise-bounds-2},
together with the Arzel\`a--Ascoli theorem (\Cref{thm:arzela-ascoli}),
show that
$\DistanceOperator{\BaseCarrier}{\MetricSymbol}\colon
\HilbertDomain_\BaseCarrier\to\ContinuousFunctions(\BaseCarrier)$
is compact.
Since the inclusion
$\ContinuousFunctions(\BaseCarrier)\to\HilbertDomain_\BaseCarrier$
is bounded,
the operator on
$\HilbertDomain_\BaseCarrier$
is compact as well.
Symmetry of the kernel and Fubini's theorem
\cite[Theorem~2.8.20(a)]{BreitGmeineder2026}
imply self-adjointness.

\paraid{p-9cf941c90fec}
If
$\DistanceOperator{\BaseCarrier}{\MetricSymbol}\HilbertFunction =\Eigenvalue  \HilbertFunction $
in
$\HilbertDomain _\BaseCarrier $
and
$\Eigenvalue \ne0$,
then
$\Eigenvalue ^{-1}\DistanceOperator{\BaseCarrier}{\MetricSymbol}\HilbertFunction $
is a continuous representative of
$\HilbertFunction $.
The representative is unique because
$\SelectedMeasure{\BaseCarrier}{\MetricSymbol}$
has full support.
\end{proof}

\paraid{p-f6f6a5a89ba1}
For every eigenfunction with nonzero eigenvalue, we use its unique continuous representative.
Apply \Cref{prop:spectral-basics} to
$\DistanceOperator{\BaseCarrier}{\MetricSymbol}$
on
$\HilbertDomain_\BaseCarrier$,
and choose an orthonormal family
$\{\Eigenfunction_\CoordinateIndex\}_{\CoordinateIndex\in\EigenIndexSet}$
with corresponding nonzero eigenvalues
$\{\Eigenvalue_\CoordinateIndex\}_{\CoordinateIndex\in\EigenIndexSet}$
as in that theorem.
For
$\SpectralCutoff >0$,
put
\[
 \EigenIndexSet _{\BaseCarrier ,\SpectralCutoff }=\{\CoordinateIndex \in \EigenIndexSet \mid|\Eigenvalue _\CoordinateIndex |>\SpectralCutoff \},
\]
and
\[
 \SpectralSubspace{\BaseCarrier }{\SpectralCutoff }
 =\LinearSpan \{\Eigenfunction _\CoordinateIndex \mid \CoordinateIndex \in \EigenIndexSet _{\BaseCarrier ,\SpectralCutoff }\}.
\]
This is a finite-dimensional space of continuous functions.
Define the orthogonal projection
$\SpectralProjection{\BaseCarrier}{\SpectralCutoff}
 \colon\HilbertDomain_\BaseCarrier\to\SpectralSubspace{\BaseCarrier}{\SpectralCutoff}$
by
\begin{equation}\label{eq:spectral-projection-definition}
 \SpectralProjection{\BaseCarrier }{\SpectralCutoff }\HilbertFunction
 =\sum_{\CoordinateIndex \in \EigenIndexSet _{\BaseCarrier ,\SpectralCutoff }}\langle \HilbertFunction ,\Eigenfunction _\CoordinateIndex \rangle_{\HilbertDomain _\BaseCarrier }\Eigenfunction _\CoordinateIndex .
\end{equation}
Equivalently,
\[
 \begin{aligned}
 \SpectralSubspace{\BaseCarrier }{\SpectralCutoff }
 =\bigoplus_{\substack{\Eigenvalue \in\spec(\DistanceOperator{\BaseCarrier}{\MetricSymbol})\\|\Eigenvalue |>\SpectralCutoff }}
       \ker(\DistanceOperator{\BaseCarrier}{\MetricSymbol}-\Eigenvalue \IdentityMap),
 \end{aligned}
\]
and
\[
 \range\SpectralProjection{\BaseCarrier }{\SpectralCutoff }=\SpectralSubspace{\BaseCarrier }{\SpectralCutoff }.
\]
These definitions do not depend on the orthonormal basis
$\{\Eigenfunction_\CoordinateIndex\}_{\CoordinateIndex\in\EigenIndexSet}$.
For every
$\HilbertFunction\in\HilbertDomain_\BaseCarrier$,
the vector
$\SpectralProjection{\BaseCarrier}{\SpectralCutoff}\HilbertFunction$
converges in the norm of
$\HilbertDomain_\BaseCarrier$
to the orthogonal projection of
$\HilbertFunction$
onto
$(\ker\DistanceOperator{\BaseCarrier}{\MetricSymbol})^\perp$
as
$\SpectralCutoff\downarrow0$.
